\subsection{Подсветка}

\begin{codebox}[led\_control]
\verb|led_control(led_id=255, r=0, g=0, b=0)| --- цвет светодиода. \verb|r|, \verb|g|, \verb|b| --- яркость от 0 до 1 (значения вне диапазона обрезаются). \verb|led_id| от 0 до 255, 255 --- все светодиоды; число вне диапазона 0--255 даёт \verb|ValueError|. Возвращает \verb|True|, если команда отправлена. Вызов без аргументов гасит всё.
\end{codebox}

В симуляторе у Мини~2 два светодиода: 0 --- левый, 1 --- правый. По исходникам 0.11.0 на настоящей плате Пионера Мини~2 подсветка одна: библиотека записывает цвет в общий модуль светодиодов, а \verb|led_id| не использует, так что любой номер красит всю подсветку. Если вам нужны разные цвета у разных светодиодов, проверьте это на своём дроне.

\begin{codefile}[python]{python-onboard/pioneer_sdk2/02_led_effects.py}
\end{codefile}

\paragraph{Что происходит.} Три эффекта на светодиодах. \verb|rainbow()|: модуль \verb|colorsys| переводит оттенок (от 0 до 1) в RGB, а сдвиг \verb|led / LED_COUNT| даёт соседним светодиодам разные цвета, и «радуга» бежит по кругу. \verb|running_light()|: каждый шаг гасит всё и зажигает следующий светодиод. \verb|pulse()|: яркость каждого канала умножается на уровень \verb|level|, так меняется яркость без смены оттенка. \verb|set_logger(False)| отключает печать ответа на каждую команду --- в анимации их сотни. В примере \verb|LED_COUNT = 4|: число светодиодов зависит от платы и версии (в симуляторе у Мини~2 их два), поэтому узнайте число на своём дроне.

Цвета смешиваются аддитивно: \verb|r=1, g=0.4| даёт жёлтый, \verb|r=1, b=1| --- фиолетовый, все каналы по 1 --- белый. Такие сочетания удобно использовать, чтобы показывать состояние программы в полёте.

\begin{taskbox}[Светофор и двоичное число]
Сделайте «светофор»: красный 3~с, жёлтый 1~с, зелёный 3~с, мигающий зелёный. Напишите функцию \verb|show_number(n)|, показывающую число от 0 до 15 в двоичном виде на четырёх светодиодах. Свяжите цвет с высотой по дальномеру: ниже 0.5~м --- зелёный, выше 1.5~м --- красный.
\expected{светофор переключается по расписанию; \texttt{show\_number(5)} зажигает 1-й и 3-й светодиоды; цвет меняется вместе с высотой.}
\end{taskbox}

\begin{rulebox}[Цвет в SDK 2]
Цвет --- доли от 0 до 1, а не числа 0..255, как в \texttt{pioneer\_sdk}. Передали \verb|r=255| --- получили полную яркость красного без ошибки: значение просто обрезается до 1. Ошибки не будет, но ожидаемого цвета тоже.
\end{rulebox}

\subsection{Груз: магнитный захват и механический захват}

В зависимости от дрона груз берёт магнит или механический захват. SDK 2 управляет обоими.

\begin{codebox}[Магнитный захват груза (Базовый Пионер)]
\verb|cargo_grab()| --- включить магнит. \verb|cargo_release()| --- выключить магнит. \verb|cargo_set(grab)| --- то же с булевым аргументом: \verb|True| включает, \verb|False| выключает. Все три возвращают \verb|True|, если успешно; если плата не поддерживает магнитный захват, печатают сообщение и возвращают \verb|False|.
\end{codebox}

\begin{codebox}[Механический захват (Пионер Мини 2)]
\verb|grab_close(movement_time=0, velocity=20)| --- закрыть захват.\\
\verb|grab_open(movement_time=0, velocity=20)| --- открыть захват.\\
\verb|grab_stop()| --- остановить движение захвата.\\
\verb|movement_time| --- сколько секунд двигать захват; 0 --- двигать до упора. \verb|velocity| --- скорость в процентах, от 0 до 100. Все три метода возвращают \verb|bool|.
\end{codebox}

\begin{codefile}[python]{python-onboard/pioneer_sdk2/17_delivery_grab.py}
\end{codefile}

\paragraph{Что происходит.} Сценарий из этапов «погрузка --- перелёт --- сброс --- возврат». На земле захват открывается (\verb|movement_time=0| --- до упора), светодиод жёлтый просит вложить груз, через 5~секунд захват закрывается со скоростью 60\%. Затем дрон взлетает, летит в точку $A$ и в точку $B$ на высоте перелёта и снижается над точкой $B$: груз сбрасывают с малой высоты, а не с большой. Каждый этап подсвечен своим цветом --- по светодиодам видно, на каком шаге программа. В \verb|finally| вызывается \verb|grab_stop()|, чтобы привод не продолжал движение. В симуляторе \verb|grab_close()| подхватывает груз под дроном так же, как магнит Базового.

\begin{infobox}[Пионер Базовый]
На Пионере Базовом вместо захвата стоит магнит: \verb|cargo_grab()|, \verb|cargo_release()| или \verb|cargo_set(True/False)|. Сценарий тот же, замените вызовы захвата.
\end{infobox}

\begin{infobox}[Важные детали захвата]
\begin{itemize}
    \item При \verb|movement_time > 0| метод ждёт столько секунд (внутри стоит \verb|time.sleep|) и только потом возвращается; при \verb|movement_time = 0| метод возвращается сразу, а привод продолжает двигаться до упора. Поэтому после \verb|grab_close()| без времени сделайте паузу, а при необходимости остановите привод \verb|grab_stop()|.
    \item Скорость \verb|velocity| вне диапазона 0--100 не вызывает исключения: библиотека печатает сообщение и возвращает \verb|False|.
    \item Скорость по умолчанию зависит от версии: в 0.16 (справочник симулятора) она 20, в 0.11.0 --- 100. Чтобы поведение не зависело от версии, всегда указывайте \verb|velocity| явно, как в примере (60 при захвате, 100 при сбросе).
    \item Если плата захват не поддерживает, методы печатают «Модуль захвата не поддерживается» и возвращают \verb|False|, а программа продолжает работу: проверяйте результат.
\end{itemize}
\end{infobox}

\begin{taskbox}[Доставка груза]
Измените пример 17: дрон берёт груз, стоящий под точкой старта, перелетает в $(1.5, 1)$ на высоту 1~м и отпускает груз, затем возвращается домой через \verb|rtl()| (с ожиданием посадки, как в разделе про возврат). Проверяйте результат каждого \verb|grab_*|; если \verb|grab_close()| вернул \verb|False|, не взлетайте.
\expected{в симуляторе груз перемещается в новую точку; при неисправном захвате программа печатает причину и не взлетает.}
\end{taskbox}

\begin{taskbox}[Захват по времени]
Закройте захват на 0.5~с, потом на 1~с с \verb|velocity| 50, замерьте \verb|time.time()| вокруг каждого вызова. Затем вызовите \verb|grab_close()| без времени и сразу \verb|grab_stop()|.
\expected{время вызова первых двух близко к 0.5 и 1~с; вызов без времени возвращается практически мгновенно, привод останавливается по \texttt{grab\_stop()}.}
\end{taskbox}

\subsection{Пульт RC}

Помимо программного управления, дроном можно управлять «как с пульта»: SDK умеет отправлять в автопилот значения каналов и читать канал пульта. Это нужно, например, чтобы написать свой пульт или передать управление с геймпада.

\begin{alertbox}[Пульт: подготовка и ограничения]
Описание SDK требует для управления каналами выставить параметры автопилота: \texttt{Copter\_man\_rcMode0=6.0}, \texttt{Copter\_man\_rcMode1=3.0}, \texttt{Copter\_man\_rcMode2=3.0}, \texttt{Copter\_flyWithoutRc=1.0}, \texttt{SensorMux\_rc=2.0}. Запомните их прежние значения и возвращайте в конце программы. Эти параметры меняют поведение автопилота, поэтому проверяйте программу на привязанном дроне, готовом к перехвату.
\end{alertbox}

\begin{codebox}[Методы RC (по 0.11.0)]
\verb|send_rc_channels(channel_1=0, channel_2=0, channel_3=0, channel_4=0,| \verb|channel_5=1, channel_6=0, channel_7=1, channel_8=0)| --- отправить значения восьми каналов; каждое число от $-1$ до $1$. Возвращает \verb|bool|. Каналы нужно отправлять \emph{постоянно}, это нужно для поддержания связи с автопилотом (частота в описании не указана).\\
\verb|get_rc_channel()| --- значение одного канала пульта (целое число) или \verb|None|. По описанию Pioneer Code это 8-й канал (тумблер SWA); какой канал читает метод, задано в настройках платы.\\
\verb|set_rc_trigger(callback)| --- вызывать функцию \verb|callback(value)|, когда значение этого канала (8-го, тумблер SWA) меняется.\\
\verb|rc_sdk1_to_sdk2(...)| --- перевод значений каналов из SDK~1 в SDK~2 (см. предупреждение ниже).
\end{codebox}

\paragraph{Каналы.} Первые четыре канала --- стики, каждый от $-1$ до $1$; каналы 5--8 --- тумблеры пульта.

\begin{center}
\small
\begin{tabularx}{\linewidth}{@{}l l X@{}}
\toprule
\textbf{Канал} & \textbf{Орган} & \textbf{Значения} \\
\midrule
\verb|channel_1| & правый стик, крен & $-1$ влево, 0 центр, $+1$ вправо \\
\verb|channel_2| & правый стик, тангаж & $-1$ вперёд, 0 центр, $+1$ назад \\
\verb|channel_3| & левый стик, газ & $+1$ вверх, 0 центр, $-1$ вниз \\
\verb|channel_4| & левый стик, рыскание & $+1$ налево, 0 центр, $-1$ направо \\
\verb|channel_5| & тумблер SWC & 0 / 1 / 2 (вверх / центр / вниз) \\
\verb|channel_6| & тумблер SWD & 0 / 2 \\
\verb|channel_7| & тумблер SWB & 0 / 1 / 2 \\
\verb|channel_8| & тумблер SWA & 0 / 2 \\
\bottomrule
\end{tabularx}
\end{center}

По умолчанию \verb|channel_5| и \verb|channel_7| равны 1, остальные --- 0. Таблица --- по сборнику примеров Pioneer Code.

\paragraph{Плавные стики.} Резкий рывок стика --- резкий манёвр. Пример отправляет каналы каждые 50~мс даже когда стики в нейтрали (автопилоту нужен непрерывный сигнал) и плавно ведёт стик по трапециевидному профилю: разгон, удержание, торможение.

\begin{center}
\begin{tikzpicture}[>=latex, xscale=1.3, yscale=1.2]
\draw[->] (0,0) -- (6.4,0) node[right] {$t$};
\draw[->] (0,-0.2) -- (0,1.4) node[above] {стик};
\draw[very thick, GeoskanBlue] (0,0) -- (2,1) -- (4,1) -- (6,0);
\draw[dashed] (2,0) -- (2,1) (4,0) -- (4,1);
\node[below, font=\small] at (1,0) {разгон};
\node[below, font=\small] at (3,0) {удержание};
\node[below, font=\small] at (5,0) {торможение};
\end{tikzpicture}
\end{center}

\begin{codefile}[python]{python-onboard/pioneer_sdk2/12_rc_ramp.py}
\end{codefile}

\paragraph{Что происходит.} Исходные параметры запоминаются до изменения и возвращаются в \verb|finally|. Функция \verb|ramp()| строит профиль из трёх участков одинаковой длительности. Запись \verb|**{axis: value}| распаковывает словарь в именованные аргументы: одна функция управляет любым стиком по его имени. Знаки: \verb|pitch| $-1$ --- вперёд, положительный \verb|yaw| --- поворот налево.

\begin{codebox}[Пример: вперёд две секунды (по 0.11.0)]
\small
\begin{verbatim}
import time
from pioneer_sdk2 import Pioneer

pioneer = Pioneer()
names = ["Copter_man_rcMode0", "Copter_man_rcMode1",
         "Copter_man_rcMode2", "Copter_flyWithoutRc", "SensorMux_rc"]
new = [6.0, 3.0, 3.0, 1.0, 2.0]
old = [pioneer.get_param(n) for n in names]
try:
    for n, v in zip(names, new):
        pioneer.set_param(n, v)
    start = time.time()
    while time.time() - start < 2:
        pioneer.send_rc_channels(channel_2=-0.3)   # вперёд
        time.sleep(0.05)
finally:
    pioneer.send_rc_channels()             # все стики в нуле
    for n, v in zip(names, old):
        if v is not None:
            pioneer.set_param(n, v)
    pioneer.close_connection()
\end{verbatim}
\end{codebox}

Пример посылает каналы 20 раз в секунду --- такой частотой SDK~1 повторял команды скорости; для SDK~2 нужная частота в описании не указана, подберите её на дроне. Значения умножаются на 32760 (каналы 1--4) и 32767 (остальные) и округляются, так что точность $1/32760$.

\paragraph{Чтение канала пульта.} \verb|get_rc_channel()| вернёт значение сразу. \verb|set_rc_trigger(callback)| запускает слежение в фоновом потоке: поток проверяет канал примерно десять раз в секунду и вызывает \verb|callback(value)|, когда значение изменилось (в том числе первый раз, когда оно прочитано). Функция выполняется в служебном потоке --- правило прежнее: она короткая, она ставит флаг, а действие делает основной цикл. Снять триггер в SDK нельзя --- он живёт до закрытия соединения; поэтому не регистрируйте триггер в цикле.

\begin{warnbox}[rc\_sdk1\_to\_sdk2 в 0.11.0 считает с ошибкой; в 0.16 не проверялось]
Задумано, видимо, перевести числа 1000--2000 из SDK~1 в $-1\ldots1$ по формуле $(x - 1500)/500$. Но в исходниках 0.11.0 скобок нет (\verb|x - 1500 / 500|), поэтому из каждого значения вычитается 3, а не пересчитывается шкала; кроме того, третий и четвёртый каналы считаются от уже изменённых первых двух. Например, \verb|rc_sdk1_to_sdk2(1500, 1500, 1500, 1500, 2000)| вернёт \verb|(1497.0, 1497.0, 1494.0, 1494.0, 1)|, а не нули. Не используйте функцию на 0.11.0; пересчитывайте шкалу сами: \verb|(x - 1500) / 500|. Пятый канал она переводит правильно: 2000 даёт 1, любое другое число --- 0.
\end{warnbox}

\begin{infobox}[RC в симуляторе]
В симуляторе RC доступен через пульт или его эмуляцию; параметры автопилота в нём не настраиваются (\verb|get_param| возвращает \verb|None|), поэтому подготовку параметров из примера в симуляторе пропускайте. Подробности зависят от версии симулятора --- проверяйте на ней.
\end{infobox}

\paragraph{Режимы программы по тумблеру.} Пока программа идёт, тумблер можно использовать как переключатель режимов:

\begin{codefile}[python]{python-onboard/pioneer_sdk2/13_rc_trigger_modes.py}
\end{codefile}

\paragraph{Что происходит.} \verb|set_rc_trigger(func)| регистрирует функцию, которая вызывается в отдельном потоке \emph{при изменении} положения тумблера SWA и получает значение 8-го канала. Обработчик только меняет глобальную переменную \verb|mode|, а основной цикл решает, что делать в каждом режиме: «событие меняет состояние --- цикл действует по состоянию». \verb|get_rc_channel()| читает текущее значение 8-го канала напрямую.

\begin{taskbox}[Тумблер как кнопка]
Добавьте третий режим, который запускает миссию из примера 05 по тумблеру. Сделайте «аварийную кнопку»: при переключении тумблера в полёте --- немедленная посадка (флаг в обработчике, \verb|land()| в основном цикле).
\expected{тумблер переключает режимы без остановки основного цикла; при аварийном переключении дрон садится из основного цикла, а не из обработчика.}
\end{taskbox}

\begin{taskbox}[Пересчёт шкалы SDK 1]
Напишите свою функцию \verb|rc1_to_rc2(ch1, ch2, ch3, ch4, ch5)|, которая переводит каналы 1000--2000 в $-1\ldots1$ по формуле $(x-1500)/500$ и пятый канал в 1 при 2000, иначе 0. Проверьте на значениях 1000, 1500, 2000.
\expected{1000 даёт $-1$, 1500 даёт $0$, 2000 даёт $1$; для пятого канала 2000 даёт 1, остальные --- 0.}
\end{taskbox}

\begin{taskbox}[Свой пульт на клавишах]
На привязанном дроне (или в симуляторе, если эмуляция пульта доступна) напишите цикл: каждые 0.05~с отправляйте \verb|send_rc_channels| с каналами, которые зависят от нажатых в консоли клавиш (вперёд/назад/влево/вправо), в конце отправьте нули. Следите, чтобы скорость была малой.
\expected{пока клавиша удерживается, дрон движется в нужную сторону; после выхода из цикла отправлены нулевые значения, дрон зависает.}
\end{taskbox}
